\newcommand{\docshorttitle}{SUPPL.\ MEMO P\&A --- ``CUSTOMER'' DOCTRINE \S 425.17(c) (CGC-25-631801)}
\newcommand{\doctitleblock}{PLAINTIFF'S SUPPLEMENTAL MEMORANDUM OF POINTS AND AUTHORITIES IN FURTHER OPPOSITION TO DEFENDANTS' SPECIAL MOTION TO STRIKE --- ``CUSTOMER'' DOCTRINE UNDER \S 425.17(c) AND THE THIRD-PARTY CLAIMANT (INSURER \& PANEL COUNSEL)}
\newcommand{\docsubblock}{(Cal.\ Rules of Court, rules 3.1113, 3.1306(c); supplemental to Plaintiff's previously-filed Opposition; cross-noticed for hearings of May 12, 2026 and May 22, 2026, Dept.\ 301)}
% Pleading-paper preamble for APRIL-24-2026 cover sheets and oral-argument
% outlines (CGC-25-631801 and CGC-25-631802). Matches house style from
% APRIL-23-SAC-DEFENSE/court-pdfs/REPLY-SAC-PREAMBLE.tex, simplified for
% single-to-three-page artifacts that do not require TOC/TOA infrastructure.
%
% Cross-hearing caption (CAPTION-801/802.tex): use \CaptionDocTitleBlock,
% \CaptionDocSubBlock, and the crosshearingsschedule environment (defined below)
% so long \doctitleblock text and hearing dates stay inside the 3in minipage.
\documentclass{article}
\usepackage[utf8]{inputenc}
\usepackage{graphicx}
\usepackage{geometry}
    \geometry{top=1in, bottom=2.2in, left=1.0in, right=1.0in, textheight=336pt}
    \setlength{\footskip}{70pt}
\usepackage{ulem}
\usepackage{fancyhdr}
    \renewcommand{\headrulewidth}{0pt}
    \renewcommand{\footrulewidth}{0.4pt}
\usepackage{parskip}
    \setlength{\parskip}{0mm}
\usepackage{quoting}
    \quotingsetup{leftmargin=0.5in, rightmargin=0.5in, vskip=-1.5mm}
    \makeatletter
        \g@addto@macro\quoting\singlespacing
        \g@addto@macro\quoting{\vspace{-2mm}}
    \makeatother
\usepackage{mathptmx}
\usepackage{amssymb}
\renewcommand{\rmdefault}{ptm}
\renewcommand{\normalsize}{\fontsize{12}{12}\selectfont}
\newcommand*{\ptm}{\fontfamily{ptm} \selectfont \fontsize{12}{12}\selectfont}
\usepackage{setspace}
    \doublespacing
\newcommand{\tab}{\hspace*{.0in}}
\newenvironment{tightcenter}{%
    \setlength\topsep{0pt}
    \setlength\parskip{0pt}
    \begin{center}
}{%
    \end{center}
}
\usepackage{tikz}
\usetikzlibrary{calc}
% Pleading-paper vertical double rule at the left margin.
\AddToHook{shipout/background}{%
    \begin{tikzpicture}[remember picture, overlay]
        \draw[line width=0.5pt] ($(current page.north west) + (0.75in,0)$) -- ($(current page.south west) + (0.75in,0)$);
        \draw[line width=0.5pt] ($(current page.north west) + (0.80in,0)$) -- ($(current page.south west) + (0.80in,0)$);
    \end{tikzpicture}%
}
\usepackage[left, pagewise]{lineno}
\setlength{\linenumbersep}{0.35in}
\renewcommand{\linenumberfont}{\normalfont\small}
\usepackage{titlesec}
\usepackage[hidelinks]{hyperref}
\usepackage{bookmark}
\usepackage{textcomp}
\usepackage{longtable}
\usepackage{booktabs}
\usepackage{array}
\usepackage{multirow}
\usepackage{calc}
% --- Cross-hearing caption cover (3in minipage): compact title + 3-column hearing table ---
% Use on first page only; keeps long \doctitleblock / \docsubblock on one sheet with party caption.
\newcommand{\CaptionDocTitleBlock}{%
  \begingroup
  \fontsize{10}{11.5}\selectfont
  \bfseries
  \raggedright
  \setlength{\parskip}{0pt}%
  \noindent\doctitleblock\par
  \endgroup
}
\newcommand{\CaptionDocSubBlock}{%
  \begingroup
  \footnotesize
  \raggedright
  \setlength{\parskip}{0pt}%
  \noindent\docsubblock\par
  \endgroup
}
% #1 = department number only (e.g.\ 301 or 302) for the line ``Dept.\ #1''.
\newenvironment{crosshearingsschedule}[1]{%
  \par\begingroup
  \footnotesize
  \setlength{\tabcolsep}{2.2pt}%
  \renewcommand{\arraystretch}{1.04}%
  \setlength{\parskip}{0pt}%
  \setlength{\parindent}{0pt}%
  \noindent\textbf{Cross-noticed for pending defense hearings (this case, Dept.\ #1):}\par\vspace{2pt}%
  \begin{tabular}{@{}>{\raggedright\arraybackslash}p{0.11\linewidth}>{\raggedright\arraybackslash}p{0.34\linewidth}>{\raggedright\arraybackslash}p{0.53\linewidth}@{}}%
  \textbf{\#} & \textbf{Motion} & \textbf{Date / dept / judge} \\ \hline
}{%
  \end{tabular}\par\endgroup\smallskip
}
% Pandoc pipe-table column widths use \real{#}.
\newcommand{\real}[1]{\pgfmathparse{#1}\pgfmathresult}
\titleformat{\section}{\fontsize{13}{13}\selectfont\normalfont\bfseries\uppercase}{}{0em}{}
\titleformat{\subsection}{\fontsize{12}{12}\selectfont\normalfont\bfseries}{}{0em}{}
\titleformat{\subsubsection}{\normalfont\bfseries}{}{0em}{}
\titlespacing*{\section}{0pt}{6pt}{2pt}
\titlespacing*{\subsection}{0pt}{4pt}{2pt}
\titlespacing*{\subsubsection}{0pt}{2pt}{2pt}
\newcommand{\calrules}[1]{Cal. Rules of Court, rule #1}
\newcommand{\ccp}[1]{Code Civ. Proc., \textsection~#1}
\newcommand{\evidcode}[1]{Evid. Code, \textsection~#1}
\providecommand{\tightlist}{%
  \setlength{\itemsep}{0pt}\setlength{\parskip}{0pt}}
% Footer: page N of total + document short-title.
\usepackage{lastpage}
\pagestyle{fancy}
\fancyhf{}
\fancyfoot[C]{\thepage{} of \pageref{LastPage} \\[1pt] \footnotesize \docshorttitle}
\fancyfoot[R]{}

\begin{document}
% Cross-hearing caption for CGC-25-631801 (Dept. 301, Hon. Van Aken).
\nolinenumbers
\begin{singlespace}
\noindent Franciscus Dylan Rosario \\
7624 Melody Drive \\
Rohnert Park, CA 94928 \\
Tel: 650.488.7510 \\
E: soltrinox@gmail.com \\
Plaintiff, In Pro Per
\end{singlespace}
\vspace*{9mm}
\begin{tightcenter}
\textbf{SUPERIOR COURT OF THE STATE OF CALIFORNIA} \\
\textbf{COUNTY OF SAN FRANCISCO}
\end{tightcenter}
\vspace*{3mm}
\begin{minipage}[t]{3in}
    \raggedright
    \textbf{FRANCISCUS DYLAN ROSARIO,} \\ \textbf{PLAINTIFF}, \\
    \hspace{1cm} v. \\
    \small
    Subhi Abdelhalim; CSAA Insurance Exchange; Phillips, Spallas \& Angstadt LLP; Priya D. Navaratnasingham; Alberto Reyna; Michael R. Chambers; Carbone, Smith \& Koyama; and Does 1 through 50, inclusive;\\
    \normalsize
    \textbf{DEFENDANTS}
    \makebox[3in]{\hrulefill}
\end{minipage}%
\hspace{3mm}%
\begin{minipage}[t]{0.5pt}
    \vspace{0pt}
    \rule{0.5pt}{2.42in}
\end{minipage}%
\hspace{3mm}%
\begin{minipage}[t]{3in}
    \raggedright
    \noindent\textbf{Case No.: CGC-25-631801}\par\smallskip
    \CaptionDocTitleBlock\smallskip
    \CaptionDocSubBlock\smallskip
    \begin{singlespace}
    \setlength{\parindent}{0pt}
    \begin{crosshearingsschedule}{301}
    (1) & MFL / leave (SAC) & \textbf{Apr.\ 30, 2026} \\
    (2) & Anti-SLAPP (\ccp{425.16}) & \textbf{May 12, 2026, 9:30 a.m.};\ Dept.\ 301;\ Hon.\ Christine Van Aken \\
    (3) & Strike (\ccp{436}) & \textbf{May 22, 2026};\ Dept.\ 301;\ Hon.\ Christine Van Aken \\
    (4) & Demurrer & Hrng.\ date \textbf{[TBD]} (reserved) \\
    \end{crosshearingsschedule}
    \noindent Complaint filed: Dec.\ 23, 2025 \\
    \noindent FAC filed: Feb.\ 25, 2026 \\
    \end{singlespace}
\end{minipage}
\vspace*{8mm}
\newpage
\linenumbers


\section*{I. INTRODUCTION AND VEHICLE}

This Supplemental Memorandum of Points and Authorities is respectfully submitted in further opposition to Defendants' special motion to strike under Code of Civil Procedure section 425.16. It is offered as a supplement to --- and not as a substitute for --- Plaintiff's previously-filed Opposition and the companion Supplemental Memorandum on the structural sequencing of \ccp{425.17(c)} filed concurrently in this case. The vehicle is \calrules{3.1113} (memorandum of points and authorities) read together with \calrules{3.1306(c)} (lodging of supplemental authority for hearings).

This memorandum is directed at one discrete question that the defense's papers attempt to manufacture: whether Plaintiff, as a third-party claimant, is a ``customer'' of CSAA's regulated insurance-claims-handling product within the meaning of \ccp{425.17(c)(2)}, and whether the panel-counsel Defendants are themselves operating within the regulated commercial channel that section 425.17(c) governs. The defense contends that, because Plaintiff is not a named insured under a CSAA policy, Plaintiff is not within the audience the carve-out protects, and that panel counsel are in any event a separate, ``litigation-only'' category outside the carve-out. Both contentions fail.

That argument collides with three independent bodies of California law: (i) the Insurance Commissioner's formal regulations at 10 California Code of Regulations section 2695.1 et seq., which expressly define ``claimant'' to include both first party and third party claimants and impose the regulatory floor on ``all persons'' and ``all claims''; (ii) Insurance Code section 790.03(h), which uses the term ``claimants'' --- without qualification --- in twelve of sixteen subdivisions and uses the disjunctive ``insureds or claimants'' to confirm that both categories are protected; and (iii) the disjunctive structure of \ccp{425.17(c)(1)} and (c)(2), which reaches the conduct of ``any person primarily engaged'' in delivering goods or services in the commercial channel --- a category that includes, as \textit{JAMS, Inc.\ v.\ Superior Court} (2016) 1 Cal.App.5th 984 confirms, the lawyers retained by a commercial actor to deliver part of that actor's regulated product.

\section*{II. ISSUES PRESENTED}

\textbf{Issue 1.} Whether a third-party claimant whose claim CSAA handled, investigated (or represented that it investigated), and ultimately denied --- on CSAA's own claim form and under its own claim number --- is a ``customer'' of CSAA's regulated insurance-claims-handling product within the meaning of Code of Civil Procedure section 425.17(c)(2).

\textbf{Issue 2.} Whether panel counsel and counsel's firm, retained by CSAA to deliver --- as part of the insurance product itself --- defense services that the policy promises to the insured, are operating within the commercial channel that \ccp{425.17(c)} regulates, such that their pre-trial communications about the insurer's claims handling are themselves within the carve-out.

\textbf{Plaintiff's answer to both:} Yes.

\section*{III. ARGUMENT}

\subsection*{A. The Insurance Commissioner's Operative Definitions Under Title 10 Treat First-Party and Third-Party Claimants as a Single Regulated Class.}

The Insurance Commissioner has, by formal rulemaking, defined the operative term ``claimant'' for purposes of California claims-handling regulation. The definition is unambiguous:

\begin{quoting}
``\,`Claimant' means a first party claimant, a third party claimant, or both and includes such claimant's designated legal representative and includes a member of the claimant's immediate family designated by the claimant.''
\end{quoting}
\begin{flushright}
\small (Cal.\ Code Regs., tit.\ 10, \S~2695.2(c).)
\end{flushright}

\noindent The companion definitions resolve any residual doubt:

\begin{quoting}
``\,`First party claimant' means any person asserting a right under an insurance policy as a named insured, other insured, or beneficiary under the terms of that insurance policy.''
\end{quoting}
\begin{flushright}
\small (Cal.\ Code Regs., tit.\ 10, \S~2695.2(s).)
\end{flushright}

\begin{quoting}
``\,`Third party claimant' means any person asserting a claim against any person under an insurance policy.''
\end{quoting}
\begin{flushright}
\small (Cal.\ Code Regs., tit.\ 10, \S~2695.2(t).)
\end{flushright}

\noindent The scope and applicability provisions are correspondingly broad:

\begin{quoting}
``The regulations which follow set forth the minimum standards for the settlement of claims and shall be applicable to the handling or settlement of \textbf{all claims}, except as specifically excluded.''
\end{quoting}
\begin{flushright}
\small (Cal.\ Code Regs., tit.\ 10, \S~2695.1(a) (emphasis added).)
\end{flushright}

\begin{quoting}
``These regulations are applicable to \textbf{all persons} and to \textbf{all insurance policies}, insurance contracts, and certificates, except as specifically excluded.''
\end{quoting}
\begin{flushright}
\small (Cal.\ Code Regs., tit.\ 10, \S~2695.1(d) (emphasis added).)
\end{flushright}

The Commissioner did not carve a third-party exemption out of the regulatory floor; he extended the floor to ``all persons'' and to ``all claims.'' The defense's ``no-privity'' theory cannot be reconciled with this text. The Commissioner has, on formal authority delegated by the Legislature under Insurance Code section 790.10, classified third-party claimants as members of the regulated class.

\subsection*{B. Insurance Code Section 790.03(h) Uses the Term ``Claimants'' Without Qualification and Treats Insureds and Claimants as Coextensive Protected Categories.}

The regulations summarized above implement Insurance Code section 790.03(h), the operative text of which is reproduced verbatim in the Supplemental Request for Judicial Notice filed concurrently in this case. Twelve of the sixteen subdivisions reach communications and conduct directed at claimants generally:

\begin{itemize}
\tightlist
\item \textbf{\S~790.03(h)(1)}: ``Misrepresenting to \textbf{claimants} pertinent facts or insurance policy provisions \dots''
\item \textbf{\S~790.03(h)(2)}: ``Failing to acknowledge and act reasonably promptly upon \textbf{communications with respect to claims} \dots''
\item \textbf{\S~790.03(h)(3)}: ``Failing to adopt and implement reasonable standards for the \textbf{prompt investigation and processing of claims} \dots''
\item \textbf{\S~790.03(h)(5)}: ``Not attempting in good faith to effectuate prompt, fair, and equitable settlements of \textbf{claims in which liability has become reasonably clear}.''
\item \textbf{\S~790.03(h)(10)}: ``\dots~\textbf{insureds or claimants} \dots''
\item \textbf{\S~790.03(h)(11)}: ``\dots~\textbf{an insured, claimant, or the physician of either} \dots''
\item \textbf{\S~790.03(h)(13)}: ``\dots~for the \textbf{denial of a claim} \dots''
\item \textbf{\S~790.03(h)(14)}: ``Directly advising \textbf{a claimant} not to obtain the services of an attorney.''
\item \textbf{\S~790.03(h)(15)}: ``Misleading \textbf{a claimant} as to the applicable statute of limitations.''
\end{itemize}

The Legislature uses the term ``claimants'' repeatedly and without qualification. Where it wished to distinguish between insureds and claimants, it did so expressly --- and in the disjunctive (``insureds or claimants''; ``an insured, claimant, or the physician of either''). The defense's reading would write that protection out of the statute. If the Legislature had wanted to limit the unfair-practices floor to insureds, it would have written ``insureds'' wherever it wrote ``claimants.'' It did not.

\subsection*{C. The California Supreme Court's Trilogy of \textit{Royal Globe} / \textit{Moradi-Shalal} / \textit{Zhang} Confirms That the Regulatory Duty Toward Third-Party Claimants Survived the Contraction of the Private Remedy.}

\textit{Royal Globe Insurance Co.\ v.\ Superior Court} (1979) 23 Cal.3d 880 held that a third-party claimant could maintain a private action for violation of section 790.03(h), expressly grounding the holding in the recognition that ``claimants'' are a category the Legislature intended to protect. \textit{Moradi-Shalal v.\ Fireman's Fund Ins.\ Companies} (1988) 46 Cal.3d 287 overruled \textit{Royal Globe} only on the implied-private-right-of-action question, and was explicit that the section ``is intended to regulate insurance industry settlement practices'' and that carriers ``remain bound by the duties the statute and its implementing regulations impose.'' (\textit{Moradi-Shalal}, 46 Cal.3d at pp.\ 304--305.) \textit{Zhang v.\ Superior Court} (2013) 57 Cal.4th 364 confirmed that \textit{Moradi-Shalal} did not bar a third-party claimant from asserting independent statutory claims premised on the same conduct. (\textit{Zhang}, 57 Cal.4th at pp.\ 380--384.)

The lesson is straightforward: the regulatory floor that defines third-party claimants as members of a single protected class with insureds remains the operative law of California. Continuing operation of the duty is also confirmed by \textit{Egan v.\ Mutual of Omaha Ins.\ Co.} (1979) 24 Cal.3d 809, 818--819, and \textit{Wilson v.\ 21st Century Ins.\ Co.} (2007) 42 Cal.4th 713, 720--721.

\subsection*{D. \ccp{425.17(c)} Uses ``Customer'' Functionally; the Audience Prong Is Triply Disjunctive.}

\begin{quoting}
``(2) The intended audience is \textbf{an actual or potential buyer or customer}, or \textbf{a person likely to repeat the statement to, or otherwise influence}, an actual or potential buyer or customer, or \textbf{the statement or conduct arose out of or within the context of a regulatory approval process, proceeding, or investigation} \dots''
\end{quoting}
\begin{flushright}
\small (\ccp{425.17(c)(2)} (emphasis added).)
\end{flushright}

The audience element is satisfied where any one of three independent branches is met. Even if the defense's narrow reading of branch (a) were somehow correct (it is not), branches (b) and (c) independently apply: a third-party claimant is unquestionably ``a person likely to \dots~influence'' the resolution of the claim, and every claim communication issued by a California-licensed insurer is conduct that ``arose out of or within the context of a regulatory \dots~investigation.'' \textit{Simpson Strong-Tie Co.\ v.\ Gore} (2010) 49 Cal.4th 12, 28--32, instructs that section 425.17(c) is read consistently with its remedial purpose. \textit{JAMS, Inc.\ v.\ Superior Court} (2016) 1 Cal.App.5th 984 and \textit{Stewart v.\ Rolling Stone LLC} (2010) 181 Cal.App.4th 664 reinforce that the inquiry under section 425.17(c) is functional, not contractual.

\subsection*{E. Panel Counsel Are Operating Within the Regulated Commercial Channel.}

\ccp{425.17(c)} is not limited to the commercial actor in name; its focus is on the cause of action and what it arises from. Where panel counsel and counsel's firm are retained by an insurer to deliver --- as part of the insurance product itself --- defense services that the policy promises to the insured, those attorneys are themselves operating within the commercial channel that \ccp{425.17(c)} regulates. (\textit{JAMS, Inc.}, 1 Cal.App.5th at pp.\ 994--995.) Where Plaintiff's claims against panel counsel arise from representations of fact about CSAA's claim handling, about the status of an investigation, about the basis for a denial, or about the existence or non-existence of factual predicates that were communicated to the court before trial, those representations were made in the course of delivering the insurance defense product to the insurer-client and, derivatively, were made about the insurer's operations within the meaning of \ccp{425.17(c)(1)}. The intended audience was, at every stage, an actual or potential decision-maker influencing the outcome of the customer's claim --- first Plaintiff as the insured-equivalent claimant, then the regulator, and ultimately the court.

\subsection*{F. Civil Code \S~47(b) and the ``Petitioning'' Gloss of \ccp{425.16(e)} Do Not Override the Carve-Out.}

Defendants rely on Civil Code \S~47(b) and on the line of cases extending the litigation privilege to communications ``in furtherance of'' judicial proceedings. As \textit{Action Apartment Ass'n v.\ City of Santa Monica} (2007) 41 Cal.4th 1232 confirms, \S~47(b) is, in many respects, coextensive with \ccp{425.16(e)} for certain deceit-based predicates. That coextensiveness cuts both ways: if the conduct does not satisfy the petitioning concept underlying \ccp{425.16(e)}, it likewise cannot be sheltered by \S~47(b) for purposes of evading \ccp{425.17(c)}. And \S~47(b) itself contains carve-outs --- including subdivision (b)(2) for communications in furtherance of intentional destruction or alteration of physical evidence, and subdivision (b)(3) for knowing concealment of insurance policies in a judicial proceeding. Business and Professions Code \S~6128(a) supplies an additional, independent prohibition: every attorney is guilty of a misdemeanor who is ``guilty of any deceit or collusion, or consents to any deceit or collusion, with intent to deceive the court or any party.''

\subsection*{G. The CSAA Denial Letter Itself Establishes Customer Status; Insurers Do Not Deny Claims to Non-Customers.}

By the act of issuing a denial --- on its own claim form, under its own claim number, signed by its own claims personnel --- CSAA acknowledged precisely the duty whose existence the defense's anti-SLAPP papers now seek to disclaim. A California-licensed insurer does not issue claim denials to non-customers. The very act of issuing a denial communicates: (i) that the recipient has presented a claim under or against an insurance policy administered by the insurer; (ii) that the insurer has, or represents that it has, conducted an investigation under its regulatory obligations; (iii) that the insurer has reached a determination disposing of the claim; and (iv) that the insurer is communicating the disposition in fulfillment of its duty under Cal.\ Code Regs., tit.\ 10, sections 2695.7(b) and 2695.7(d). Each is, by the insurer's own conduct, an acknowledgment that the recipient is within the regulated class. Panel counsel cannot, by labeling pre-trial commercial representations as ``advocacy,'' convert that admission into anti-SLAPP-protected petitioning activity.

\subsection*{H. The Privilege of California Licensure: A Statutorily Created Duty Cannot Be Disclaimed by Re-Characterization.}

The right to operate as an insurance exchange in California is conferred in exchange for the obligation to handle every claim, first-party and third-party alike, in accordance with the regulatory floor. (Ins.\ Code, \S~700 et seq.) A carrier that wishes to disclaim third-party claims-handling duties is, in substance, asking the State for a unilateral modification of the bargain on which its certificate of authority was issued. That request is foreclosed by the statutes and regulations themselves. The same point reaches panel counsel: counsel retained by an insurer to deliver, as part of the insurance product, defense services cannot, by labeling themselves as ``litigation actors,'' step outside the channel for purposes of \ccp{425.17(c)}.

\section*{IV. APPLICATION TO CGC-25-631801}

The moving Defendants include CSAA and panel counsel retained by CSAA. The same analysis applies as to CSAA: Plaintiff is a third-party claimant within the regulated class as defined by Cal.\ Code Regs., tit.\ 10, \S~2695.2(c) and (t); CSAA's pre-litigation conduct was made in the course of delivering the regulated product; the audience element is satisfied under all three branches of \ccp{425.17(c)(2)}. As to panel counsel, the analysis is identical: counsel's pre-trial communications about the existence and adequacy of CSAA's claim investigation were communications ``made in the course of delivering'' the insurance product within the meaning of \ccp{425.17(c)(1)}; the audience for those communications --- Plaintiff, the regulator, and ultimately the court --- falls within branches (b) and (c) of \ccp{425.17(c)(2)}. The special motion to strike must be denied at the gateway as to both CSAA and panel counsel, and the litigation privilege under Civil Code \S~47(b) supplies no independent route around the carve-out where, as alleged here, the conduct is alleged to constitute extrinsic fraud aimed at impairing Plaintiff's ability to present his case at trial.

\section*{V. CONCLUSION}

The defense's effort to evade Code of Civil Procedure section 425.17(c) by re-characterizing Plaintiff as a non-customer of CSAA's regulated insurance-claims-handling product, and by re-characterizing panel counsel as ``litigation-only'' actors outside the regulated commercial channel, cannot be reconciled with the text and structure of California insurance law. Plaintiff respectfully requests that the Court deny the special motion to strike at the gateway under \ccp{425.17(c)} as to both CSAA and the panel-counsel Defendants, without reaching the merits-stage probability-of-prevailing inquiry, with the consequence that no interlocutory appeal lies under \ccp{425.17(e)} and the ordinary civil docket may proceed.

\input{SIGBLOCK.tex}
\end{document}
